\chapter{Phase Conclusions and Future Work}
\section{Verified Primary Results}
\subsection{Public action policies provide stable execution baselines}
The official \pifive{} weights completed 117/120 episodes across four LIBERO tasks. Success rate dropped significantly when images were zeroed out or frozen, indicating that this reproduction indeed relies on visual input; executing old trajectories or geometric scripts cannot substitute this result. This part falls under public method reproduction and input intervention, not a self-trained general action model with equivalent capability.
\subsection{The Benefits of Chinese Understanding Need to Be Distinguished from Hierarchical Control}
Direct Chinese input achieves 22/30, frozen hierarchical version achieves 28/30, and only translation handed to flat policy achieves 29/30. Hierarchical planning does not outperform the only-translation control; subsequent experiments with fixed switching and frozen visual inspection also did not show stable additional gains. Correct high-level descriptions may still fail to alter the low-level policy's target preference; this interface issue is more direct than continuing to scale up the language model.
\subsection{The checker improves offline evaluation but does not improve overall task performance}
The second version of the checker improves completion recall on approximate completion examples, but single false alarms still exceed the preset threshold. In exploratory closed-loop comparisons, both the baseline and adapted versions achieved 28/30; the adapted version, however, confirmed correct completions less frequently. Two consecutive votes can mitigate some sporadic errors but cannot guarantee elimination of relevant misjudgments. Milestone memory improves stop decisions on small batches of new initial states within a single scenario, but remains insufficient to demonstrate cross-skill generalizable completion judgment.
\subsection{Explicit operation points produce measurable effects during local execution}
Visual affordance low-rank adaptation significantly improves drawer handle localization; strict hit rate on the same input reaches 93.7\% in the 2B version. After depth back-projection and planning, positions achieve 4/4 for LoRA, 0/4 for base, and 4/4 for ground-truth localization. Here, validation is for reaching, not full grasping. Subsequently, without resetting state, handoff to the action policy yields only 2/4 for the full task; native policy baseline reveals that additional pre-action steps may disrupt the stage distribution familiar to the policy.
\subsection{Edge deployment has completed accuracy comparison, but speed is still unsatisfactory}
The restored Jetson supports GPU four-bit inference and adapter loading capabilities, and has completed a PC-to-edge comparison with 350 frames and 1014 queries. The 2B quantized version achieved a strict hit rate of 91.1\%, but text output per query matched exactly only 76.33\%. Three edge-side visual-to-PC pre-contact joint tests passed; under the 15W mode, with 512-pixel input, the median latency was 5.062 seconds and the 95th percentile was 7.149 seconds, failing to meet the single-query 3-second target. Reducing image size resulted in limited latency variation, and the new size still lacks full precision regression.
\section{Correction of Previous Results}
The explicitly targeted adaptation branch later discovered an implementation error in the low-rank weight merging formula. Among six saved inputs, the old deployment differed from the unmerged reference by approximately 0.6055 in maximum action, which was reduced to approximately 0.01144 after correction. The correction demonstrates improved numerical interface performance, but no full closed-loop retest has been conducted with the corrected version. Therefore, the old low score cannot be directly attributed to the ineffectiveness of this type of adaptation method, nor can the success rate be assumed to increase after correction.

The real robotic arm successfully moved objects along a fixed trajectory, but repeated attempts resulted in collisions with the table. It should serve as data and experience for the execution interface, not as evidence of reliable autonomous grasping. Three SmolVLA demonstrations and restricted trials did not establish stable task behaviors, and 50\% of actions were not actually completed. Simulation results are not used to replace these real hardware conclusions.
\section{Engineering Experience Gained from Multiple Stages}
\begin{enumerate}
\item First verify the input and action semantics, then explain the model's capabilities. Any error in image rotation, wrist bias, normalization, action blocks, or low-rank merging may obscure the true cause.
\item Keep simple baselines for the new modules. Translation, fixed switching, training set waypoints, native policy, and ground-truth localization each answer different questions and cannot be compared solely by the final success rate of a single complex pipeline.
\item Make evaluation consistent with usage. The checker should see starting frames, current frames, and prompts consistent with closed-loop operation; accuracy should be re-evaluated when generation length, quantization range, and input size change.
\item Hierarchically validate localization, planning, reaching, handover, task success, and real-time performance. A passing of a prior layer can help diagnose failures in the subsequent layer, but does not imply the entire pipeline has passed.
\item Retain failed episodes and version conditions. Conditions such as different initial states, reruns, test set participation, and ideal depth will alter the interpretation of the numbers.
\end{enumerate}
\section{Future Work Ranked by Risk}
\subsection{Restoring Trustworthy Action Adaptation Baselines}
First, perform a small-scale numerical consistency check on the branch after weight merging, then redo the target interface comparison for frozen initial states. Only after confirming that the deployed actions are consistent with the training interface should the training data be expanded or the adaptation scheme be re-evaluated.
\subsection{Extended Visual Affordance and Stage Compatibility}
Maintain native policy reference, validate operation point localization and stage identification separately; prioritize testing camera changes, occlusions, and new layouts to clarify when planning pre-contact is needed versus when existing skills should be directly applied. Cross-scenario, real depth error, and multiple skills remain in the scope for future work.
\subsection{On-device Optimization: First Identify Bottlenecks, Then Modify Runtime}
Profiling points to many small operators and four-bit linear-computation overhead. Future optimization should measure output length, power mode, and runtime separately. Every change must retain matched-input accuracy checks, peak-memory measurements, and swap records. Pausing simulation while waiting for the service does not establish real-time compliance.
\subsection{Expanding Independent Validation Rather Than Simply Increasing Demonstrations}
Current key comparisons have only 3--5 episodes, making it difficult to reliably estimate tail risks. In the next round, tasks, initial states, budgets, checkpoint selection methods, and failure categorization should be pre-fixed, while cross-task validation should be increased; debugging states should be separated from formal testing. Real hardware validation requires separate recording of collision margins, calibration, and stopping conditions, to be conducted only after successful simulation.
